\section{问题重述与数据分析}

\subsection{题目要求}
本题要求实现可复用的财务问答 Agent\cite{financeqacourse}。给定一家公司和一道英文问题，系统要从财务数据库和年报中找证据，自主决定何时查询、检索或计算，并给出有依据的回答。课程以 Snowflake 为例，问了两件事：“公司的收入增长总体呈现怎样的趋势？”“公司重点强调了哪些产品与技术举措？”回答第一题要对齐财年、比较收入并计算增速；回答第二题要从年报中找出具体的产品和技术举措。我们要验证，同一套系统能否按题目切换查证方式，并在换到新的公司后独立完成问答。表~\ref{tab:fq-question-types} 展示了不同问题所需的证据。

\begin{table}[ht]
\centering
\caption{问题类型与所需证据}
\label{tab:fq-question-types}
\begin{tabular}{@{}p{.28\columnwidth}p{.63\columnwidth}@{}}
\toprule
问题类型 & 核心作答要求\\
\midrule
指标趋势与比较 & 找到可比期间、定义一致的指标，说明单位、方向和变化幅度\\
产品与技术举措 & 定位公司年报中的业务或研发陈述，概括行动及对应来源\\
综合财务分析 & 联合查询数值与年报解释，区分已发生结果和管理层计划\\
\bottomrule
\end{tabular}
\end{table}

\subsection{数据组成与实验目标}
数据包含 6 家公司、9 份年报和 SQLite 财务数据库，另有 56 道问题。开发集覆盖 Snowflake、Moderna、ServiceNow 和 Norfolk Southern 的 37 道题；测试集覆盖未参与开发的 Palantir 和 Abbott，各 10 题和 9 题。

实验分两步检验：先看增强检索能否在开发集提高支持率，再看选定的方案能否答完两家新公司的问题，以及在哪类题上仍会出错。

\section{Agent 设计与实现}
如图~\ref{fig:fq-workflow} 所示，Agent 每轮接收原问题和当前工具结果，再决定继续查询或生成答案\cite{yao2023react}。它能够查看当前公司的年报目录、搜索和展开相邻年报片段、查找财务指标标签、提取年度事实记录、执行只读 SQL 及计算同比变化。每题最多 6 轮，每轮最多提出 3 个工具请求。检索片段保留公司、文件名、报告日期、章节和表格上下文；事实记录保留指标标签、期间和单位。答案附带来源编号，程序以该题实际收到的工具结果验证编号。

年报解析后去除脚本、样式及隐藏 XBRL 标记，按正文段落和表格分块，并使用 SQLite FTS5 建立全文索引。处理的是课程提供的年报，不读取具体问题或答案。公司范围与 SQL 数据库均以题目中的 CIK 限定。年度财务事实区分有明确年度标记的流量、期末余额和期间尚未核实的记录；检索到多条披露时不直接求和，并保留出处以便复核。

\begin{figure*}[t]
\centering
\includegraphics[width=.93\textwidth]{financeqa-workflow.pdf}
\caption{Agent 的决策循环。Flash 可以发出工具请求、读回结果后继续查证，也可以直接提交答案和来源编号。}
\label{fig:fq-workflow}
\end{figure*}

\section{实验设置}
基础方案使用 BM25 对年报片段排序\cite{robertson2009bm25}，按概念词匹配财务标签；增强方案扩展常见财务同义词，并按原查询词命中、短语命中、检索分数及报告日期重排结果，对指标标签加入通用概念映射，并在期间可靠性相同时优先当前年度申报。两种方案均由 DeepSeek Flash 作答，提示词和轮数相同，温度为 0，每轮生成上限为 1600 token。以开发集支持题数选定方案，再冻结代码、提示词、模型配置和数据索引。测试集全部答案保存后才运行课程评测。

每道题的独立检查器将回答与原题和参考答案对照，给出“支持”“矛盾”或“信息不足”。报告分别统计这三类数量；总体支持率为支持题数除以题目数。按课程判定规则，提供部分支持信息也可归入“支持”。评分使用同一 DeepSeek V4 Pro 模型和课程原始检查提示。
